\begin{claim}{D.61}{Proved}{Full generation and exact negative conjugator equations}
Put $F=\widehat F(x,y)$ and $c=[x,y]$.
Every closed subgroup $J\leq F$ satisfying
$c\in\overline{[J,J]}$ is $F$ itself.
Consequently every continuous endomorphism of $F$
carrying $c$ to a conjugate of $c^\mu$, with
$\mu\in\widehat\Z^\times$, is an automorphism.

Let $b$ be any negative full once-punctured-torus
outer involution, with the specified extension
\[
 E=F\rtimes\langle v\rangle,\quad
 v^2=1,\quad vcv^{-1}=c^{-1}.
\]
Let $A_2$ be its action on $H_1(F,\Z_2)$.
Set $v_i=c^iv$ and
$\alpha_i=\operatorname{Inn}(v_i)|_F$, for $i=0,1$.
On a separate free profinite group
$F_0=\widehat F(X,Y)$, with $c_0=[X,Y]$, use the
two standard actual negative involutions
\[
 \begin{array}{c|cc}
      &X&Y\\ \hline
 \beta_0&YXY^{-1}&Y^{-1}\\
 \beta_1&Y&X .
 \end{array}                                             \tag{D.61a}
\]
For $j=0,1$, let $\mathcal T_j(b)$ be the space
of peripheral outer isomorphisms $[\phi]:F_0\to F$
conjugating $[\beta_j]$ to $b$.
The peripheral requirement is
$\phi(c_0)\sim_F c^\mu$ for a full unit $\mu$.
Identifying $F_0$ with $F$ by $X\mapsto x$,
$Y\mapsto y$ identifies this with the usual
outer conjugator space.

For $A_2\bmod2=I$, define
\[
 \begin{split}
 \mathcal S_{0,i}=\{(\mu,h,w):\
 &\mu\in\widehat\Z^\times,\ h\in F,\quad
 w^2=1,\ w\in E\setminus F,\\
 &w\text{ belongs to the other full section class},\\
 &[h,w]=1,\quad [h,v_i]=c^\mu\}.
 \end{split}                                             \tag{D.61b}
\]
Here \emph{other} means the class distinct from the
common cusp-reflection class of D.58.
There is a continuous bijection, with continuous
inverse,
\[
 \mathcal S_{0,0}\sqcup\mathcal S_{0,1}
       \xrightarrow{\ \cong\ }\mathcal T_0(b),\qquad
 (\mu,h,w)\longmapsto
             [\,X\mapsto h,\ Y\mapsto v_iw\,].            \tag{D.61c}
\]
For $A_2\bmod2\ne I$, define
\[
 \mathcal S_{1,i}
   =\{(\mu,f)\in\widehat\Z^\times\times F:
                         [f,\alpha_i(f)]=c^\mu\}.
                                                               \tag{D.61d}
\]
The corresponding continuous bijection is
\[
 \mathcal S_{1,0}\sqcup\mathcal S_{1,1}
       \xrightarrow{\ \cong\ }\mathcal T_1(b),\qquad
 (\mu,f)\longmapsto
             [\,X\mapsto f,\ Y\mapsto\alpha_i(f)\,].
                                                               \tag{D.61e}
\]
Equivalently, writing $z=fv_i$, the last equation
is
\[
                z\in E\setminus F,\qquad
                         [z^2,v_i]=c^\mu.               \tag{D.61f}
\]
The two basis images are then $zv_i,v_iz$.
Thus no separate generation condition occurs in
either parameter space. The equations and their
outer conjugator spaces are identified exactly;
nonemptiness for an arbitrary $b$ remains unproved.
\end{claim}

\begin{proof}
\emph{The boundary test for a full closed subgroup.}
Suppose $J$ were proper. A proper open subgroup
$U\leq F$ contains $J$: separate a point outside
the compact closed subgroup $J$ in a finite
quotient and take the inverse image of the image
of $J$. In particular $c\in U$.
Write $\Pi=F(x,y)$ for the dense discrete
punctured-torus group. Its subgroup
$\Pi_U=\Pi\cap U$ has the same finite index
$d=[F:U]>1$, and its full profinite completion
is exactly $U$.

In the associated connected covering surface,
the boundary lift through the base point has
degree one because $c\in\Pi_U$.
The degrees of all boundary lifts sum to $d$.
There must therefore be a second boundary
component. In the homology of an orientable
surface with at least two boundary components,
the boundary classes have only their total-sum
relation. Each individual boundary class is
nonzero over every $\F_p$.
The chosen class is the class of $c$, with its
specified base path. Its image in
$H_1(U,\F_p)=H_1(\Pi_U,\F_p)$ is nonzero.
On the other hand
$c\in\overline{[J,J]}\subseteq\overline{[U,U]}$,
so that image is zero. This contradiction proves
$J=F$.

Now let $\rho:F\to F$ be an endomorphism with
$\rho(c)=q c^\mu q^{-1}$ for a full unit $\mu$.
Its image is compact and closed. In the conjugate
image $J=q^{-1}\rho(F)q$, the element $c^\mu$
belongs to $\overline{[J,J]}$.
Taking its $\mu^{-1}$ power puts $c$ in the same
closed commutator subgroup. The result just proved
makes $\rho$ onto.
A finitely generated profinite group is Hopfian:
the inverse-image operation of a surjective
endomorphism permutes the finite set of open
normal subgroups of each given index, so its
kernel lies in every such subgroup. Their
intersection is trivial. Thus $\rho$ is an
automorphism.
The same argument applies to a homomorphism
$F_0\to F$ after the explicit basis identification
in the claim.

\emph{The exact two standard extensions.}
The substitutions in (D.61a) square to the
identity and send $c_0$ to $c_0^{-1}$.
Let $E_j=F_0\rtimes_{\beta_j}\langle v_0\rangle$.
For $j=0$ put $w_0=Y^{-1}v_0$.
There are mutually inverse continuous maps
\[
 \begin{aligned}
 (\widehat\Z\langle X\rangle\times\langle w_0\rangle)
                   \amalg\langle v_0\rangle
       &\longrightarrow E_0,\\
 (X,w_0,v_0)&\longmapsto(X,Y^{-1}v_0,v_0),\\
 E_0&\longrightarrow
  (\widehat\Z\langle X\rangle\times C_2)\amalg C_2,\\
 (X,Y,v_0)&\longmapsto(X,v_0w_0,v_0).
 \end{aligned}                                             \tag{D.61g}
\]
Both displayed finite generators have order
two. The equations $w_0^2=1$ and $[X,w_0]=1$
follow directly from the action $\beta_0$.
Conversely $Y=v_0w_0$ gives
$v_0Yv_0=Y^{-1}$ and
$v_0Xv_0=YXY^{-1}$.
Substitution on the displayed generators proves
that the two maps are inverse.
The quotient label sends $X$ to zero and
$w_0,v_0$ to one in $C_2$.
The two finite factors represent distinct
section classes: the character taking $w_0$
to one and $X,v_0$ to zero separates them.

For $j=1$ put $z_0=Xv_0$. The exact maps are
\[
 \begin{aligned}
 \langle v_0\rangle\amalg\widehat\Z\langle z_0\rangle
       &\longrightarrow E_1,\\
 (v_0,z_0)&\longmapsto(v_0,Xv_0),\\
 E_1&\longrightarrow C_2\amalg\widehat\Z,\\
 (X,Y,v_0)&\longmapsto(z_0v_0,v_0z_0,v_0).
 \end{aligned}                                             \tag{D.61h}
\]
Here $v_0$ swaps the two displayed images of
$X,Y$, so all semidirect relations hold.
The two composites fix the generators.
The quotient label is one on both $v_0$ and
$z_0$. These maps specify the full profinite
products through their universal properties.

\emph{Normalization of every peripheral outer conjugator.}
Take $[\phi]\in\mathcal T_j(b)$.
After an inner change in $F$, choose a
representative satisfying
\[
                         \phi(c_0)=c^\mu .              \tag{D.61i}
\]
The unit $\mu$ is unique: a conjugator between
two full unit powers of $c$ normalizes
$L=\overline{\langle c\rangle}$; the full
peripheral normalizer inside $F$ is $L$ itself,
which centralizes $c$.

The centre-free extension construction gives
the unique isomorphism
$\Phi:E_j\to E$ restricting to $\phi$ and
inducing the identity on $C_2$.
Explicitly, conjugation by $\phi$ carries
the full inverse image of
$\langle[\beta_j]\rangle$ in
$\operatorname{Aut}(F_0)$ to the full inverse
image of $\langle b\rangle$ in
$\operatorname{Aut}(F)$.
The embeddings of these extensions are faithful:
an element centralizing their free kernel has
trivial quotient label and then belongs to
the trivial centre of that kernel.
This gives both the extension map and its
inverse.

Put $t=\Phi(v_0)$. It has order two,
has nonzero quotient label, and inverts
$c^\mu$ by (D.61i).
As $\mu$ is a unit, it also inverts $c$.
D.58's actual cusp normalizer therefore gives
the unique expression
\[
                           t=c^zv,\qquad z\in\widehat\Z.
                                                               \tag{D.61j}
\]
Let $i\in\{0,1\}$ be the residue of $z$ modulo
two and put $r=(i-z)/2\in\widehat\Z$.
Replacing $\Phi$ by $\operatorname{Inn}(c^r)\Phi$
retains (D.61i) and sends (D.61j) to
$c^{z+2r}v=v_i$.
Thus every outer class has a representative
with the exact two normalization equations
\[
                    \phi(c_0)=c^\mu,\qquad
                              \Phi(v_0)=v_i.            \tag{D.61k}
\]

This representative and the index $i$ are
unique within the outer class.
Indeed two representatives satisfying the
first equation differ by an inner conjugation
by an element of $C_F(c^\mu)=L$.
Conjugation by $c^r$ changes the reflection
exponent by $2r$. Equality of two indices in
$\{0,1\}$ modulo two first makes the indices
equal, and then $2r=0$ gives $r=0$.
All peripheral centralizers and normalizers
used here are the full specified subgroups.

\emph{The two-class parameter map and its inverse.}
Suppose first that $A_2\bmod2=I$.
For a tuple in (D.61b), put $k=v_iw\in F$.
Commutation of $h$ with $w$ gives the exact
identity
\[
                       [h,k]=[h,v_i]=c^\mu.             \tag{D.61l}
\]
The homomorphism $X\mapsto h$, $Y\mapsto k$
is consequently an automorphism $F_0\to F$
by the boundary test.
Also
\[
 \alpha_i(k)=k^{-1},\qquad
                         \alpha_i(h)=khk^{-1},           \tag{D.61m}
\]
the second equation following by inserting
$w h w=h$.
Equations (D.61m) prove
$\phi\beta_0=\alpha_i\phi$.
Since $\alpha_i$ represents $b$, this is
the asserted peripheral outer conjugator.
Its actual extension sends $v_0$ to $v_i$
and $w_0$ to $w$.

Conversely, normalize any member of
$\mathcal T_0(b)$ by (D.61k).
Put $h=\phi(X)$ and $w=\Phi(w_0)$.
The source relations in (D.61g) give
$w^2=1$, $w\notin F$, and $[h,w]=1$.
That isomorphism carries the two distinct
source section classes to the two distinct
target classes, so $w$ is in the other class.
Since $\phi(Y)=v_iw$, equation (D.61i)
gives (D.61l), hence (D.61b).
These two constructions recover the exact
tuple and the normalized representative.
This proves bijectivity in (D.61c).

\emph{The one-class parameter map and its inverse.}
Now let $A_2\bmod2\ne I$.
A tuple in (D.61d) defines
$\phi(X)=f$ and $\phi(Y)=\alpha_i(f)$.
Its commutator is $c^\mu$, so the boundary
test again proves full invertibility.
The equality $\alpha_i^2=1$ makes
$\phi\beta_1=\alpha_i\phi$.
This gives the indicated outer conjugator
and an extension sending $v_0$ to $v_i$.
Conversely, for the normalized representative
of any member of $\mathcal T_1(b)$, the same
intertwining equation forces
$\phi(Y)=\alpha_i(\phi(X))$.
Taking $f=\phi(X)$ and using (D.61i) yields
exactly (D.61d).

Finally $z=fv_i$ gives
$z^2=f\alpha_i(f)$ and
\[
 [z^2,v_i]
   =f\alpha_i(f)\bigl(\alpha_i(f)f\bigr)^{-1}
   =[f,\alpha_i(f)] .                                   \tag{D.61n}
\]
Conversely $f=zv_i$ and
$\alpha_i(f)=v_iz$.
This proves the alternative parameterization
and both basis-image formulas in (D.61f).

\emph{Topology and the remaining existence obligation.}
Each parameter space is closed in the indicated
compact product: the full other section class
in (D.61b) is the compact image of its conjugation
map, and all equations are continuous.
The maps to automorphisms are continuous on
generator images, and hence so are their maps
to the peripheral outer isomorphism space.
The latter space is Hausdorff and closed in
the profinite outer group, after the explicit
basis identification. The continuous
bijections just proved are therefore
homeomorphisms.

The two integral homology actions in (D.61a)
have respectively identity and nonidentity
reduction modulo two. Conjugacy preserves
this distinction, so only the indicated
standard action can occur in each row.
The formulas supply full actual generators
for every tuple in the displayed spaces.
They establish no tuple for an arbitrary
negative full outer involution.
The remaining problem is the existence of
solutions of those precise equations,
with their retained peripheral element and
reflection index.
The executable
\texttt{code/part\_c/torus\_negative\_parameters\_d61.py}
checks the exact standard free-word and
free-product identities and the boundary
test on all labelled connected covers of
degrees two through five.
\end{proof}
